\section{Forcing in every perfect normal strategy}\label{sec:forcing}
We work after Lemma~\ref{lem:support}, in faithful normal standard form. The following precise base-game facts are the imports from \cite[Lemmas 2.1 and 2.2, and Section 2]{TypeGames}.

\begin{lemma}[Checked classes and the base bridge]\label{lem:base}
A perfectly matched coarse partition is a sharp centralizer partition, with its unique standard commutant mirror. Fine effects whose marginals are sharp are the corresponding joint spectral projections. In a perfect $\GIII$ strategy, the flag effects $F_j$ form a centralizer partition of one. Each copy has a centralizer $M_4$ on $1-F_j$ and a partial isometry $w_j$ satisfying
\begin{equation}\label{eq:bridge}
 w_j^*w_j=F_j,\quad w_jw_j^*=C_j:=F_{j+1}+F_{j+2},
 \quad \sigma_t^\phi(w_j)=2^{it}w_j.
\end{equation}
In a copy write $f=F_j$, choose a minimal core projection $q$ and centralizer maps $u_x$ from $q$ to $c_x$. If $D_x$ is the $T_x$ sign difference, its pullback is
\[
 X_x=(f+u_x^*)D_x(f+u_x)\in(f+q)M(f+q).
\]
The positive $(T_x,U_0)$ bridge cell gives $X_xH\xi=Y_0H\xi$, where $H=f+q/\sqrt2$ and $Y_0$ is the pullback of the original Bob sign difference. The extracted bridge is a core translate of
\begin{equation}\label{eq:beta}
 \beta=q(f+u_x^*)D_x(f+u_x)f.
\end{equation}
\end{lemma}
Here the identity involving $H$ is the positive-cell calculation in the proof of the cited modular-bridge lemma. It uses calibrated sign supports and rejected sign events, before any assumption of sharp individual $T_x$ signs. Thus the same calculation applies to a calibrated variant of that question. We spell out its consequence because it is essential to the labelling argument.

For the checked-class assertion, let $E_c,F_d$ be matching coarse effects on the two sides. Zero mismatched events give $E_cF_d\xi=0$ for $c\ne d$. Summation gives $E_c\xi=F_c\xi$. It follows that $E_cE_d\xi=0$ for $c\ne d$, and separation makes the partition sharp. A self-adjoint commutant mirror implies $\phi(E_cz)=\phi(zE_c)$, so it is centralizer and its mirror is $JE_cJ$. Positive fine effects lie below their sharp marginals; summing them identifies their joint spectral projections.

The base square then gives a centralizer $M_4$ with arbitrary multiplicity. Its matrix units belong to the algebra generated by its rays. For the bridge, calibration fixes sign support $f+c_x$. Pullback yields commuting left and right contractions on the standard $(f+q)$ corner. The four parity cells have joint-support vectors $(f+(-1)^{xy}q/\sqrt2)\xi$. The three positive cells identify the common reflection; the negative cell makes it anticommute with $f-q$. Its off-diagonal part is \eqref{eq:beta}, with initial projection $f$, final projection $q$ and modular eigenvalue two. The base flag links identify its designated range with $C_j$, giving \eqref{eq:bridge}. These are exactly the bridge conclusions used below.

\subsection{Common unlabelled marginals}
Every occurrence of a particular label matches the same binary Bob question. Lemma~\ref{lem:base} therefore makes it one sharp centralizer involution, independent of its Alice context. A fine joint effect lies below its sharp label projection. Consequently its unlabelled marginal commutes with every label in that context.

For a labelled row question, its comparisons to the three original Bob columns fix the flag and all three cell-sign marginals. Those are the original sharp marginals. The checked-class assertion identifies every fine unlabelled row effect with the original joint spectral projection. Hence every labelled row has its original unlabelled POVM.

Consider a labelled $T_x$ in context $t$. Calibration against the common original Bob $C_1$ fixes the residual effects and the sum of the two sign effects, $f+c_x$. Write their difference as $D_x^t$ and pull it back:
\[
 X_x^t=(f+u_x^*)D_x^t(f+u_x).
\]
The inherited positive $(T_x,U_0)$ cell is present for this variant and the original question. The calibrated core, the original Bob question and the vector $H\xi$ are the same in both calculations. Thus the positive-cell identity of Lemma~\ref{lem:base} gives
\[
 X_x^tH\xi=Y_0H\xi=X_xH\xi.
\]
On the standard corner $e=f+q$, the centralizer element $H$ is invertible. The vector $H\xi$ separates $eMe$. Subtraction therefore gives $X_x^t=X_x$. Since $f+u_x$ is a partial isometry from $e$ onto $f+c_x$, pushing forward recovers $D_x^t=D_x$. The fixed sign sum and residual effects identify the entire unlabelled $T_x$ POVM. This proves common unlabelled marginals for all nine variants of every base question.

\subsection{Commutation with the bridge and the block equations}
Each of the twelve labels occurs in a context attached to every base question. The common-marginal result consequently makes each label commute with every original Alice effect in its copy. It commutes with the core matrix units, with $D_x$ and with \eqref{eq:beta}, hence with its core translates $w_j$. This conclusion does not assert commutation of incompatible labels with one another.

The auxiliary square relations now hold as operator identities on the full identity. Their internal magic-square calculation gives $X_jZ_j=-Z_jX_j$. For example, writing the first two rows as $a,b,ab$ and $c,d,cd$, the first two columns give $ac,bd$, and the last row and column give $acbd=-abcd$. Together with the row/column commutations these imply $cb=-bc$ and $ad=-da$. Here $a=X_j,d=Z_j$.

The three $(\mathrm{Lab},U_q)$ questions make $U_q$ centralizer involutions commuting with every $F_j$. To justify conditional contexts without assuming sharp original effects, let their on-block effects be $E_s$. Matching the block bit to $\mathrm{Lab}$ gives $\sum_sE_s=F_D$. If a matched centralizer sign projection $P_t$ is incompatible with $s$, the losing event gives
\[
 E_sJP_tJ\xi=0.
\]
Since $JP_tJ\xi=P_t\xi$, separation yields $E_sP_t=0$. Thus each $E_s$ lies under all its compatible sign projections. Summing gives $[F_D,P_t]=0$ and identifies the restricted joint spectral projections. Missing tuples impose exactly their stated relation on $F_D$. There is no conclusion off that block.

In particular $X_j,Z_j$ commute with $C_j$ by their commutation with $w_j$, and the two $Z$ contexts give
\begin{equation}\label{eq:swap}
 Z_jC_j=F_{j+1}-F_{j+2},\qquad X_jF_{j+2}X_j=F_{j+1}.
\end{equation}
The remaining contexts give
\begin{equation}\label{eq:Yblocks}
\begin{aligned}
 Y_q^jF_j&=U_qF_j,\\
 Y_b^jC_j&=F_{j+1}U_a+F_{j+2}U_b,\\
 Y_c^jC_j&=F_{j+1}U_b+F_{j+2}U_a,\\
 Y_a^jC_j&=X_jU_cC_j,
 \qquad [X_j,U_q]C_j=0\quad(q=a,b,c).
\end{aligned}
\end{equation}
For $q=c$ the commutation is furnished by the active commuting triple; the other two are the explicit joint contexts.

\subsection{The extracted Cuntz pair}
Put $p=j+1$, $r=j+2$ and
\begin{equation}\label{eq:extracted}
 e_{j,0}=F_pw_j,\qquad e_{j,1}=X_jF_rw_j,\qquad
 S_x=\sum_j e_{j,x}^*.
\end{equation}
The bridge and flag-swap identities imply
\[
 e_{j,0}e_{j,0}^*=e_{j,1}e_{j,1}^*=F_p,\qquad
 e_{j,0}e_{j,1}^*=0,\qquad
 \sum_x e_{j,x}^*e_{j,x}=F_j.
\]
For distinct $j$ their initial supports $F_j$ and their common final supports $F_{j+1}$ are disjoint. Expanding the finite sums therefore gives
\[
 S_x^*S_y=\delta_{xy}\un,\qquad
 S_0S_0^*+S_1S_1^*=\un.
\]
Every $S_x$ has modular eigenvalue $1/2$, since the flags and labels are centralizer and the $w_j$ have eigenvalue two.

Commuting $Y_b^j,Y_c^j$ through $w_j$, and using \eqref{eq:Yblocks}, gives
\[
 e_{j,x}U_q=U_{q|x}e_{j,x}\quad(q=b,c),
\]
where $b|0=a$, $b|1=b$, $c|0=b$, $c|1=a$. For $x=1$, the imposed $X_j/U_a$ and $X_j/U_b$ commutations move the section observable to the left. For the active generator the same transport, the flag swap and $Y_a^j=X_jU_c$ on $C_j$ give
\[
 e_{j,0}U_a=U_ce_{j,1},\qquad e_{j,1}U_a=U_ce_{j,0}.
\]
Taking adjoints and summing in $j$ gives all six equations \eqref{eq:six} for the global $U_q$ and $S_x$. Theorem~\ref{thm:A} applies inside the faithful support host.

\begin{remark}[The limit of label commutation]\label{rem:literal}
Requiring every label to commute with every already-labelled fine effect would be inconsistent: summing the fine effects of a context containing $Z_j$ would make $X_j$ commute with $Z_j$. Their anticommutation forbids this. The property proved and used is commutation with the unlabelled base marginals and the extracted bridge. It permits all the incompatible square contexts.
\end{remark}
